% --- Archon physics formalization source begin ---
% archon:physics
% archon:covers PhyXMiniProblems/problem_phyx_mini_0674.lean
% archon:source-report reports/phyx_mini/problem_phyx_mini_0674.source.json

\chapter{Physics problem phyx\_mini\_0674}
\label{ch:PhyXMiniProblems_problem_phyx_mini_0674}

\paragraph{Problem source.}
\#\# Physical scenario

A blue car of length \textbackslash{}( 4.52 \textbackslash{}text\{ m\} \textbackslash{}) is moving north on a roadway that intersects another perpendicular roadway as shown in figure. The width of the intersection from near edge to far edge is \textbackslash{}( 28.0 \textbackslash{}text\{ m\} \textbackslash{}). The blue car has a constant acceleration of magnitude \textbackslash{}( 2.10 \textbackslash{}text\{ m/s\}\textasciicircum{}2 \textbackslash{}) directed south. The time interval required for the nose of the blue car to move from the near (south) edge of the intersection to the north edge of the intersection is \textbackslash{}( 3.10 \textbackslash{}text\{ s\} \textbackslash{}).

\#\# Question

How far is the nose of the blue car from the south edge of the intersection when it stops?

\#\# Answer choices

- A: A: 18.3m

- B: B: 34.5m

- C: C: 35.9m

- D: D: 27.6m

\#\# Figure caption (auxiliary; use the image as primary evidence)

The image shows a top-down view of a four-way intersection. There are two cars in the scene:

1. **Red Car**: 
   - Located on the left side, moving horizontally towards the right.
   - Accompanied by a vector labeled \textbackslash{}(a\_R\textbackslash{}) indicating its acceleration direction towards the right.

2. **Blue Car**: 
   - Positioned at the bottom, moving vertically upwards.
   - Displayed with a vector labeled \textbackslash{}(v\_B\textbackslash{}) indicating its velocity direction upwards.
   - Another vector labeled \textbackslash{}(a\_B\textbackslash{}) shows its acceleration direction downwards.

The roads have dashed yellow lines marking the lanes. There is a compass rose on the upper right side indicating directions with labels:
- N for North (upwards)
- S for South (downwards)
- E for East (rightwards)
- W for West (leftwards)

A measurement line, labeled "28.0 m," is drawn on the right side, indicating the width of the intersection.

\#\# Dataset metadata

Kinematics; Physical Model Grounding Reasoning, Multi-Formula Reasoning

\paragraph{Recorded answer/context.}
C: C: 35.9m

\paragraph{Figure/image path.}
/root/proposal\_for\_physic/science-mango/phyx\_mini\_run/phyx\_data/test\_image/674.png

\paragraph{Formalization target.}
create a compiling Lean file with sorry bodies at `PhyXMiniProblems/problem\_phyx\_mini\_0674.lean`. The Lean declarations must preserve the physical quantities, dimensions or dimensional roles, figure labels, governing-law hypotheses, and final relation expressed by this problem.
Use Mathlib/Physlib names found through LeanExplore where available. If a physics API is missing, introduce faithful local abstractions rather than scalar placeholder aliases.

\begin{theorem}[Physics formalization target]
\label{thm:physics:phyx_mini_0674:target}
\lean{PhyXMiniProblems.ProblemPhyXMini0674.problem_phyx_mini_0674}
\uses{def:physics:phyx-mini-0674:phyxminiproblems-problemphyxmini0674-roadintersectionsetup, def:physics:phyx-mini-0674:phyxminiproblems-problemphyxmini0674-matchesproblemstatement, def:physics:phyx-mini-0674:phyxminiproblems-problemphyxmini0674-matchesprimaryfigure, def:physics:phyx-mini-0674:phyxminiproblems-problemphyxmini0674-hasphysicalstoppingevent, def:physics:phyx-mini-0674:phyxminiproblems-problemphyxmini0674-satisfiesconstantaccelerationkinematics, def:physics:phyx-mini-0674:phyxminiproblems-problemphyxmini0674-stoppingdistanceinmeters, def:physics:phyx-mini-0674:phyxminiproblems-problemphyxmini0674-displayedstoppingdistancemeters, def:physics:phyx-mini-0674:phyxminiproblems-problemphyxmini0674-isclosestdisplayedstoppingdistance}
The assigned autoformalize agent should translate this physics problem into Lean declarations in the covered file, with theorem and lemma proof bodies written as `by sorry`.
\end{theorem}
\begin{proof}
This is an autoformalization task, not a proof task. Produce statements that can later be proved by the physics prover without weakening the physical model.
\end{proof}

% --- Archon declaration topology begin ---
\section{Declaration topology}
\begin{definition}[Length Quantity]
  \label{def:physics:phyx-mini-0674:phyxminiproblems-problemphyxmini0674-lengthquantity}
  \lean{PhyXMiniProblems.ProblemPhyXMini0674.LengthQuantity}
  A physical length or signed longitudinal position
\end{definition}
\begin{proof}
  This is the declared model definition of Length Quantity.
\end{proof}
\begin{definition}[Time Quantity]
  \label{def:physics:phyx-mini-0674:phyxminiproblems-problemphyxmini0674-timequantity}
  \lean{PhyXMiniProblems.ProblemPhyXMini0674.TimeQuantity}
  A physical time coordinate or time interval
\end{definition}
\begin{proof}
  This is the declared model definition of Time Quantity.
\end{proof}
\begin{definition}[Signed Velocity Quantity]
  \label{def:physics:phyx-mini-0674:phyxminiproblems-problemphyxmini0674-signedvelocityquantity}
  \lean{PhyXMiniProblems.ProblemPhyXMini0674.SignedVelocityQuantity}
  A signed one-dimensional velocity, with north represented by positive readout
\end{definition}
\begin{proof}
  This is the declared model definition of Signed Velocity Quantity.
\end{proof}
\begin{definition}[Signed Acceleration Quantity]
  \label{def:physics:phyx-mini-0674:phyxminiproblems-problemphyxmini0674-signedaccelerationquantity}
  \lean{PhyXMiniProblems.ProblemPhyXMini0674.SignedAccelerationQuantity}
  A signed one-dimensional acceleration
\end{definition}
\begin{proof}
  This is the declared model definition of Signed Acceleration Quantity.
\end{proof}
\begin{definition}[length Readout]
  \label{def:physics:phyx-mini-0674:phyxminiproblems-problemphyxmini0674-lengthreadout}
  \lean{PhyXMiniProblems.ProblemPhyXMini0674.lengthReadout}
  \uses{def:physics:phyx-mini-0674:phyxminiproblems-problemphyxmini0674-lengthquantity}
  Read a physical length in the length unit of a coherent unit choice
\end{definition}
\begin{proof}
  This is the declared model definition of length Readout.
\end{proof}
\begin{definition}[time Readout]
  \label{def:physics:phyx-mini-0674:phyxminiproblems-problemphyxmini0674-timereadout}
  \lean{PhyXMiniProblems.ProblemPhyXMini0674.timeReadout}
  \uses{def:physics:phyx-mini-0674:phyxminiproblems-problemphyxmini0674-timequantity}
  Read a physical time in the time unit of a coherent unit choice
\end{definition}
\begin{proof}
  This is the declared model definition of time Readout.
\end{proof}
\begin{definition}[velocity Readout]
  \label{def:physics:phyx-mini-0674:phyxminiproblems-problemphyxmini0674-velocityreadout}
  \lean{PhyXMiniProblems.ProblemPhyXMini0674.velocityReadout}
  \uses{def:physics:phyx-mini-0674:phyxminiproblems-problemphyxmini0674-signedvelocityquantity}
  Read a signed velocity in the speed unit induced by a coherent unit choice
\end{definition}
\begin{proof}
  This is the declared model definition of velocity Readout.
\end{proof}
\begin{definition}[acceleration Readout]
  \label{def:physics:phyx-mini-0674:phyxminiproblems-problemphyxmini0674-accelerationreadout}
  \lean{PhyXMiniProblems.ProblemPhyXMini0674.accelerationReadout}
  \uses{def:physics:phyx-mini-0674:phyxminiproblems-problemphyxmini0674-signedaccelerationquantity}
  Read a signed acceleration in the acceleration unit induced by a unit choice
\end{definition}
\begin{proof}
  This is the declared model definition of acceleration Readout.
\end{proof}
\begin{definition}[Car Color]
  \label{def:physics:phyx-mini-0674:phyxminiproblems-problemphyxmini0674-carcolor}
  \lean{PhyXMiniProblems.ProblemPhyXMini0674.CarColor}
  The two cars distinguished by color in the supplied figure
\end{definition}
\begin{proof}
  This is the declared model definition of Car Color.
\end{proof}
\begin{definition}[Cardinal Direction]
  \label{def:physics:phyx-mini-0674:phyxminiproblems-problemphyxmini0674-cardinaldirection}
  \lean{PhyXMiniProblems.ProblemPhyXMini0674.CardinalDirection}
  Cardinal directions printed on the compass rose
\end{definition}
\begin{proof}
  This is the declared model definition of Cardinal Direction.
\end{proof}
\begin{definition}[Compass Label]
  \label{def:physics:phyx-mini-0674:phyxminiproblems-problemphyxmini0674-compasslabel}
  \lean{PhyXMiniProblems.ProblemPhyXMini0674.CompassLabel}
  Literal compass labels shown in the image
\end{definition}
\begin{proof}
  This is the declared model definition of Compass Label.
\end{proof}
\begin{definition}[Roadway Axis]
  \label{def:physics:phyx-mini-0674:phyxminiproblems-problemphyxmini0674-roadwayaxis}
  \lean{PhyXMiniProblems.ProblemPhyXMini0674.RoadwayAxis}
  The two perpendicular roadway axes visible in the image
\end{definition}
\begin{proof}
  This is the declared model definition of Roadway Axis.
\end{proof}
\begin{definition}[Figure Placement]
  \label{def:physics:phyx-mini-0674:phyxminiproblems-problemphyxmini0674-figureplacement}
  \lean{PhyXMiniProblems.ProblemPhyXMini0674.FigurePlacement}
  Approach lanes in which the two cars are drawn
\end{definition}
\begin{proof}
  This is the declared model definition of Figure Placement.
\end{proof}
\begin{definition}[Intersection Edge]
  \label{def:physics:phyx-mini-0674:phyxminiproblems-problemphyxmini0674-intersectionedge}
  \lean{PhyXMiniProblems.ProblemPhyXMini0674.IntersectionEdge}
  The near and far intersection edges relevant to the blue car
\end{definition}
\begin{proof}
  This is the declared model definition of Intersection Edge.
\end{proof}
\begin{definition}[Figure Vector Label]
  \label{def:physics:phyx-mini-0674:phyxminiproblems-problemphyxmini0674-figurevectorlabel}
  \lean{PhyXMiniProblems.ProblemPhyXMini0674.FigureVectorLabel}
  Literal vector labels printed beside the arrows in the supplied figure
\end{definition}
\begin{proof}
  This is the declared model definition of Figure Vector Label.
\end{proof}
\begin{definition}[Intersection Figure]
  \label{def:physics:phyx-mini-0674:phyxminiproblems-problemphyxmini0674-intersectionfigure}
  \lean{PhyXMiniProblems.ProblemPhyXMini0674.IntersectionFigure}
  \uses{def:physics:phyx-mini-0674:phyxminiproblems-problemphyxmini0674-carcolor, def:physics:phyx-mini-0674:phyxminiproblems-problemphyxmini0674-cardinaldirection, def:physics:phyx-mini-0674:phyxminiproblems-problemphyxmini0674-compasslabel, def:physics:phyx-mini-0674:phyxminiproblems-problemphyxmini0674-roadwayaxis, def:physics:phyx-mini-0674:phyxminiproblems-problemphyxmini0674-figureplacement, def:physics:phyx-mini-0674:phyxminiproblems-problemphyxmini0674-intersectionedge, def:physics:phyx-mini-0674:phyxminiproblems-problemphyxmini0674-figurevectorlabel}
  Qualitative geometry and labels transcribed from the primary bitmap
\end{definition}
\begin{proof}
  This is the declared model definition of Intersection Figure.
\end{proof}
\begin{definition}[Road Intersection Setup]
  \label{def:physics:phyx-mini-0674:phyxminiproblems-problemphyxmini0674-roadintersectionsetup}
  \lean{PhyXMiniProblems.ProblemPhyXMini0674.RoadIntersectionSetup}
  \uses{def:physics:phyx-mini-0674:phyxminiproblems-problemphyxmini0674-lengthquantity, def:physics:phyx-mini-0674:phyxminiproblems-problemphyxmini0674-timequantity, def:physics:phyx-mini-0674:phyxminiproblems-problemphyxmini0674-signedvelocityquantity, def:physics:phyx-mini-0674:phyxminiproblems-problemphyxmini0674-signedaccelerationquantity, def:physics:phyx-mini-0674:phyxminiproblems-problemphyxmini0674-intersectionfigure}
  The physical quantities of the blue-car trajectory and the supplied figure. The stopping time and stopping position are unconstrained fields here; their physical relations occur only in the premise predicates below
\end{definition}
\begin{proof}
  This is the declared model definition of Road Intersection Setup.
\end{proof}
\begin{definition}[Matches Problem Statement]
  \label{def:physics:phyx-mini-0674:phyxminiproblems-problemphyxmini0674-matchesproblemstatement}
  \lean{PhyXMiniProblems.ProblemPhyXMini0674.MatchesProblemStatement}
  \uses{def:physics:phyx-mini-0674:phyxminiproblems-problemphyxmini0674-lengthreadout, def:physics:phyx-mini-0674:phyxminiproblems-problemphyxmini0674-timereadout, def:physics:phyx-mini-0674:phyxminiproblems-problemphyxmini0674-accelerationreadout, def:physics:phyx-mini-0674:phyxminiproblems-problemphyxmini0674-roadintersectionsetup}
  Numerical data stated in the prose, read in SI units. The acceleration is signed because north is positive, so the stated southward 2.10 m/s² is represented by -2.10
\end{definition}
\begin{proof}
  This is the declared model definition of Matches Problem Statement.
\end{proof}
\begin{definition}[Matches Primary Figure]
  \label{def:physics:phyx-mini-0674:phyxminiproblems-problemphyxmini0674-matchesprimaryfigure}
  \lean{PhyXMiniProblems.ProblemPhyXMini0674.MatchesPrimaryFigure}
  \uses{def:physics:phyx-mini-0674:phyxminiproblems-problemphyxmini0674-lengthreadout, def:physics:phyx-mini-0674:phyxminiproblems-problemphyxmini0674-roadintersectionsetup}
  Primary-image readout: the blue car is on the south approach and points north, the red car is on the west approach, a\_R points east, v\_B points north, a\_B points south, and the 28.0 m marker spans the south and north edges
\end{definition}
\begin{proof}
  This is the declared model definition of Matches Primary Figure.
\end{proof}
\begin{definition}[Has Physical Stopping Event]
  \label{def:physics:phyx-mini-0674:phyxminiproblems-problemphyxmini0674-hasphysicalstoppingevent}
  \lean{PhyXMiniProblems.ProblemPhyXMini0674.HasPhysicalStoppingEvent}
  \uses{def:physics:phyx-mini-0674:phyxminiproblems-problemphyxmini0674-lengthreadout, def:physics:phyx-mini-0674:phyxminiproblems-problemphyxmini0674-timereadout, def:physics:phyx-mini-0674:phyxminiproblems-problemphyxmini0674-velocityreadout, def:physics:phyx-mini-0674:phyxminiproblems-problemphyxmini0674-accelerationreadout, def:physics:phyx-mini-0674:phyxminiproblems-problemphyxmini0674-roadintersectionsetup}
  Positivity, direction, and the defining readout of the selected future stop. These conditions identify the physical branch but do not constrain the stopping position or mention a displayed answer
\end{definition}
\begin{proof}
  This is the declared model definition of Has Physical Stopping Event.
\end{proof}
\begin{definition}[Satisfies Constant Acceleration Kinematics]
  \label{def:physics:phyx-mini-0674:phyxminiproblems-problemphyxmini0674-satisfiesconstantaccelerationkinematics}
  \lean{PhyXMiniProblems.ProblemPhyXMini0674.SatisfiesConstantAccelerationKinematics}
  \uses{def:physics:phyx-mini-0674:phyxminiproblems-problemphyxmini0674-lengthreadout, def:physics:phyx-mini-0674:phyxminiproblems-problemphyxmini0674-timereadout, def:physics:phyx-mini-0674:phyxminiproblems-problemphyxmini0674-velocityreadout, def:physics:phyx-mini-0674:phyxminiproblems-problemphyxmini0674-accelerationreadout, def:physics:phyx-mini-0674:phyxminiproblems-problemphyxmini0674-roadintersectionsetup}
  One-dimensional constant-acceleration kinematics, stated relative to the instant when the blue car's nose reaches the south edge. The same relations hold in every coherent choice of units. This interface contains no stopping distance, answer value, or answer label
\end{definition}
\begin{proof}
  This is the declared model definition of Satisfies Constant Acceleration Kinematics.
\end{proof}
\begin{definition}[stopping Distance In Meters]
  \label{def:physics:phyx-mini-0674:phyxminiproblems-problemphyxmini0674-stoppingdistanceinmeters}
  \lean{PhyXMiniProblems.ProblemPhyXMini0674.stoppingDistanceInMeters}
  \uses{def:physics:phyx-mini-0674:phyxminiproblems-problemphyxmini0674-lengthreadout, def:physics:phyx-mini-0674:phyxminiproblems-problemphyxmini0674-roadintersectionsetup}
  The requested signed northward distance of the stopped nose from the south edge
\end{definition}
\begin{proof}
  This is the declared model definition of stopping Distance In Meters.
\end{proof}
\begin{lemma}[initial Velocity from edge Crossing]
  \label{lem:physics:phyx-mini-0674:phyxminiproblems-problemphyxmini0674-initialvelocity-from-edgecrossing}
  \lean{PhyXMiniProblems.ProblemPhyXMini0674.initialVelocity_from_edgeCrossing}
  \uses{def:physics:phyx-mini-0674:phyxminiproblems-problemphyxmini0674-velocityreadout, def:physics:phyx-mini-0674:phyxminiproblems-problemphyxmini0674-roadintersectionsetup, def:physics:phyx-mini-0674:phyxminiproblems-problemphyxmini0674-matchesproblemstatement, def:physics:phyx-mini-0674:phyxminiproblems-problemphyxmini0674-satisfiesconstantaccelerationkinematics}
  The crossing data determine the blue car's speed as its nose enters the intersection. This is the first of the two constant-acceleration calculations used by the problem
\end{lemma}
\begin{proof}
  The conclusion follows from the governing relations and figure readouts named in the statement of initial Velocity from edge Crossing.
\end{proof}
\begin{lemma}[stopping Distance exact]
  \label{lem:physics:phyx-mini-0674:phyxminiproblems-problemphyxmini0674-stoppingdistance-exact}
  \lean{PhyXMiniProblems.ProblemPhyXMini0674.stoppingDistance_exact}
  \uses{def:physics:phyx-mini-0674:phyxminiproblems-problemphyxmini0674-roadintersectionsetup, def:physics:phyx-mini-0674:phyxminiproblems-problemphyxmini0674-matchesproblemstatement, def:physics:phyx-mini-0674:phyxminiproblems-problemphyxmini0674-hasphysicalstoppingevent, def:physics:phyx-mini-0674:phyxminiproblems-problemphyxmini0674-satisfiesconstantaccelerationkinematics, def:physics:phyx-mini-0674:phyxminiproblems-problemphyxmini0674-stoppingdistanceinmeters}
  Using the future zero-velocity event gives the exact SI stopping distance before rounding to a displayed answer
\end{lemma}
\begin{proof}
  The conclusion follows from the governing relations and figure readouts named in the statement of stopping Distance exact.
\end{proof}
\begin{definition}[Answer Choice]
  \label{def:physics:phyx-mini-0674:phyxminiproblems-problemphyxmini0674-answerchoice}
  \lean{PhyXMiniProblems.ProblemPhyXMini0674.AnswerChoice}
  Labels of the four stopping-distance choices printed in the problem
\end{definition}
\begin{proof}
  This is the declared model definition of Answer Choice.
\end{proof}
\begin{definition}[displayed Stopping Distance Meters]
  \label{def:physics:phyx-mini-0674:phyxminiproblems-problemphyxmini0674-displayedstoppingdistancemeters}
  \lean{PhyXMiniProblems.ProblemPhyXMini0674.displayedStoppingDistanceMeters}
  \uses{def:physics:phyx-mini-0674:phyxminiproblems-problemphyxmini0674-answerchoice}
  The distance in metres printed beside an answer label
\end{definition}
\begin{proof}
  This is the declared model definition of displayed Stopping Distance Meters.
\end{proof}
\begin{definition}[Is Closest Displayed Stopping Distance]
  \label{def:physics:phyx-mini-0674:phyxminiproblems-problemphyxmini0674-isclosestdisplayedstoppingdistance}
  \lean{PhyXMiniProblems.ProblemPhyXMini0674.IsClosestDisplayedStoppingDistance}
  \uses{def:physics:phyx-mini-0674:phyxminiproblems-problemphyxmini0674-roadintersectionsetup, def:physics:phyx-mini-0674:phyxminiproblems-problemphyxmini0674-stoppingdistanceinmeters, def:physics:phyx-mini-0674:phyxminiproblems-problemphyxmini0674-answerchoice, def:physics:phyx-mini-0674:phyxminiproblems-problemphyxmini0674-displayedstoppingdistancemeters}
  A choice is closest to the stopping distance among all displayed values
\end{definition}
\begin{proof}
  This is the declared model definition of Is Closest Displayed Stopping Distance.
\end{proof}
% --- Archon declaration topology end ---
% --- Archon physics formalization source end ---
